\subsection{Ergebnisse der Umsetzung}
Durch die Umsetzung haben die Studierenden einen lokalen Cluster, in dem sie ihre Implementierungen in einer produktionsechten Umgebung ausprobieren können.
Die Automation der Cluster-Einrichtung bringt nicht nur den Vorteil, dass die Studierenden den Cluster leicht auf ihren Geräten einrichten können, ohne Kubernetes-Kenntnisse zu benötigen, sondern der Cluster muss dadurch auch nur temporär auf diesen Geräten vorhanden sein, wenn er gerade benötigt wird. 
Dadurch können die Studierenden beim Entwickeln den Cluster mit nur einem Befehl starten, ihre Implementierungen testen und nach Abschluss der Arbeit diesen wieder mit einem Cleanup-Befehl entfernen, bis sie den lokalen Cluster wieder benötigen.

Durch die Ressourcenverwaltung können nun Pods besser auf den Nodes verteilt werden, zusätzlich wird auch verhindert das ein Service alle Ressourcen für sich selbst beanspruchen kann.
Zusammen mit den Probes und den RollingUpdates erzielt das Codetask bei einem Update eine Zero Downtime.

Durch den Ingress sind nun drei Services nach außen erreichbar mit http, während die anderen Services nur im Cluster selbst erreichbar bleiben.